
\documentclass[11pt,a4paper]{article}

\usepackage[margin=1.7cm]{geometry}
\usepackage{amsmath,amssymb,mathtools,bm}
\usepackage{booktabs,longtable,array,tabularx}
\usepackage{pdflscape}
\usepackage{xcolor}
\usepackage{colortbl}
\usepackage{enumitem}
\usepackage{siunitx}
\usepackage[hidelinks]{hyperref}
\usepackage[T1]{fontenc}
\usepackage{lmodern}
\usepackage{microtype}

\definecolor{softblue}{RGB}{235,244,252}
\definecolor{softgreen}{RGB}{237,248,240}
\definecolor{softgray}{RGB}{245,245,245}
\definecolor{softred}{RGB}{252,239,239}

\newcommand{\U}{\bm{U}}
\newcommand{\Ud}{\bm{U}^{\mathrm{DEM}}}
\newcommand{\Usrc}{\bm{U}^{\mathrm{src}}}
\newcommand{\epsp}{\varepsilon}
\newcommand{\ind}{\mathbb{I}}
\newcommand{\norm}[1]{\left\lVert #1\right\rVert}
\newcommand{\dd}{\,\mathrm{d}}
\newcommand{\VSMALL}{\varepsilon_{\mathrm{num}}}

\newcolumntype{Y}{>{\raggedright\arraybackslash}X}
\newcolumntype{L}[1]{>{\raggedright\arraybackslash}p{#1}}

\title{\textbf{Mathematical Comparison of the Two Solid-Velocity Interpolation Models}\\
\large \texttt{laplaceAnchored} versus \texttt{laplaceSetValues}}
\author{Derived directly from \texttt{lambdaDotModel.C}}
\date{}

\begin{document}
\maketitle

\section{Purpose and notation}

The code first converts discrete DEM particle velocities into a cell-centred raw velocity field
$\Ud$, and then constructs an interpolated Eulerian solid velocity field $\U$.
Both vector equations are solved component by component for $U_x$, $U_y$, and $U_z$.

For cell $i$:
\begin{align}
\mathcal{P}_i
&= \{\text{DEM particles whose centres lie in cell }i\},\\
n_i
&= |\mathcal{P}_i|,\\
\Ud_i
&=
\begin{cases}
\dfrac{1}{n_i}\displaystyle\sum_{p\in\mathcal{P}_i}\bm{v}_p,
& n_i>0.5,\\[1.2ex]
\bm{0}, & n_i\le 0.5.
\end{cases}
\label{eq:rawdem}
\end{align}

The symbols used below are:
\begin{center}
\begin{tabularx}{0.96\textwidth}{L{2.5cm}Y}
\toprule
Symbol & Meaning\\
\midrule
$\epsp_i$ & \texttt{porosityF} in cell $i$\\
$\epsp_c$ & \texttt{solidPorosityCutoff}\\
$V_i$ & volume of cell $i$\\
$o_i$ & occupancy indicator, $o_i=\ind(n_i>0.5)$\\
$f=(i,j)$ & face shared by neighbouring cells $i$ and $j$\\
$A_f$ & face area\\
$d_{ij}$ & centre-to-centre distance normal to face $f$\\
$D$ & Laplace diffusivity in \texttt{laplaceSetValues}; in the code $D=10^{-4}\,\mathrm{m^2}$\\
$c_{\mathrm{DEM}}$ & strong anchor coefficient; default $10^6$\\
$c_{\mathrm{bg}}$ & weak background anchor coefficient; default $10^{-12}$\\
\bottomrule
\end{tabularx}
\end{center}

\section{Step-by-step mathematical comparison}

\begin{landscape}
\small
\setlength{\tabcolsep}{5pt}
\renewcommand{\arraystretch}{1.38}

\begin{longtable}{L{2.0cm}L{10.3cm}L{10.3cm}}
\toprule
\textbf{Step}
&
\cellcolor{softblue}\textbf{\texttt{laplaceAnchored}}
&
\cellcolor{softgreen}\textbf{\texttt{laplaceSetValues}}
\\
\midrule
\endfirsthead

\toprule
\textbf{Step}
&
\cellcolor{softblue}\textbf{\texttt{laplaceAnchored}}
&
\cellcolor{softgreen}\textbf{\texttt{laplaceSetValues}}
\\
\midrule
\endhead

1. Initial field
&
The solution field is initialized by
\[
\U_i^{(0)}=\Ud_i.
\]
Cells classified as non-solid are immediately reset to zero.
&
The solution field is initialized by
\[
\U_i^{(0)}=\Ud_i.
\]
No explicit gas-cell reset is applied before the Laplace solve.
\\

2. Occupied-cell indicator
&
\[
o_i=
\begin{cases}
1,&n_i>0.5,\\
0,&n_i\le0.5.
\end{cases}
\]
Thus a stationary particle with $\Ud_i=\bm 0$ is still recognized as occupied.
&
The hard-reference set is determined from velocity magnitude:
\[
\mathcal A
=
\left\{i:\norm{\Ud_i}>0\right\}.
\]
Therefore an occupied cell whose mean DEM velocity is exactly zero is \emph{not} included in the hard-reference set.
\\

3. Solid-domain indicator
&
\[
s_i=
\ind\!\left(\epsp_i<\epsp_c\ \lor\ o_i=1\right).
\]
An occupied cell is treated as solid even when its porosity is above the cutoff.
&
\[
\chi_i=\operatorname{pos}(1-\epsp_i).
\]
Hence the Laplace region is inferred from whether $\epsp_i$ is below unity, rather than directly from $n_i$.
\\

4. Characteristic cell length
&
\[
\ell_i^2
=
\max\!\left(V_i^{2/3},\VSMALL\right).
\]
This scales the anchor coefficient with mesh size.
&
No cell-length scaling is used in the main Laplace equation. The diffusivity $D$ is spatially constant.
\\

5. Anchor coefficient
&
First define
\[
c_i=
\begin{cases}
c_{\mathrm{DEM}},&o_i=1\ \text{or}\ s_i=0,\\
c_{\mathrm{bg}},&s_i=1\ \text{and}\ o_i=0.
\end{cases}
\]
Then
\[
a_i=\frac{c_i}{\ell_i^2}.
\]
Thus occupied and gas cells receive a large reaction coefficient, while empty solid cells receive an almost-zero coefficient.
&
There is no reaction/penalty coefficient. Instead, cells in $\mathcal A$ are imposed through an algebraic hard constraint:
\[
\U_i=\Ud_i,\qquad i\in\mathcal A.
\]
\\

6. Source field
&
Because empty DEM cells already have zero velocity,
\[
\Usrc_i=
\begin{cases}
\Ud_i,&o_i=1,\\
\bm 0,&o_i=0.
\end{cases}
\]
Gas cells are also assigned $\Usrc_i=\bm 0$.
&
The main equation has no volumetric source:
\[
\bm S_i=\bm 0.
\]
The only prescribed nonzero information enters through \texttt{setValues}.
\\

7. Face diffusion mask
&
For an internal face $f=(i,j)$, the code first interpolates $s$ to the face and then explicitly sets the coefficient to zero whenever either adjacent cell is non-solid. For the binary mask this is equivalent to
\[
\gamma_f=s_i s_j.
\]
Therefore
\[
\gamma_f=
\begin{cases}
1,&s_i=s_j=1,\\
0,&\text{otherwise}.
\end{cases}
\]
&
The cell mask is interpolated:
\[
\chi_f=\mathcal I_f(\chi).
\]
An interface face is identified by
\[
0<\chi_f<1.
\]
For those faces the matrix coupling is manually removed. This is equivalent to an effective face diffusivity
\[
D_f^{\mathrm{eff}}=
\begin{cases}
0,&0<\chi_f<1,\\
D,&\text{otherwise}.
\end{cases}
\]
\\

8. Continuous governing equation
&
The code solves a screened-Laplace or reaction-diffusion equation:
\[
-\nabla\cdot\left(\gamma\nabla\U\right)
+a\U
=
a\Usrc.
\tag{A}
\]
&
Away from hard-reference cells, the code solves the homogeneous Laplace equation:
\[
\nabla\cdot\left(D\nabla\U\right)=\bm 0.
\tag{B}
\]
Since the right-hand side is zero, reversing the overall matrix sign does not change the mathematical solution.
\\

9. Finite-volume face weight
&
For a coupled solid-solid face $f=(i,j)$:
\[
w_{ij}^{A}
=
\gamma_f\frac{A_f}{d_{ij}}.
\]
If either cell is non-solid, $w_{ij}^{A}=0$.
&
For a face that is not cut:
\[
w_{ij}^{B}
=
D\frac{A_f}{d_{ij}}.
\]
For a solid-gas interface face, $w_{ij}^{B}=0$.
\\

10. Cell equation before constraints
&
Ignoring explicit non-orthogonal corrections, cell $i$ satisfies
\[
\left(
a_iV_i+\sum_{j\in N(i)}w_{ij}^{A}
\right)\U_i
-
\sum_{j\in N(i)}w_{ij}^{A}\U_j
=
a_iV_i\Usrc_i.
\tag{A-FV}
\]
&
For an unconstrained cell:
\[
\left(
\sum_{j\in N(i)}w_{ij}^{B}
\right)\U_i
-
\sum_{j\in N(i)}w_{ij}^{B}\U_j
=
\bm 0.
\tag{B-FV}
\]
\\

11. Treatment of DEM cells
&
A DEM cell is a \emph{soft} anchor during the linear solve:
\[
a_i\gg 1
\quad\Longrightarrow\quad
\U_i\approx\Ud_i.
\]
The equality is not algebraically exact inside the matrix solve.
&
A DEM cell in $\mathcal A$ is a \emph{hard} anchor. Its matrix row is replaced by
\[
\U_i=\Ud_i.
\]
Thus the prescribed value is exact during the main Laplace solve.
\\

12. Treatment of empty solid cells
&
For $s_i=1$ and $o_i=0$:
\[
a_i=\frac{c_{\mathrm{bg}}}{\ell_i^2}\approx0.
\]
Therefore
\[
-\nabla\cdot(\gamma\nabla\U)\approx\bm0,
\]
so the field is nearly harmonic, with a very weak tendency toward zero.
&
There is no background reaction term:
\[
\nabla\cdot(D\nabla\U)=\bm0.
\]
The value is a purely harmonic extension of the hard-reference values and boundary information.
\\

13. Treatment of gas cells
&
For $s_i=0$:
\[
\gamma_f=0
\quad\text{on adjacent internal faces},
\qquad
a_i=\frac{c_{\mathrm{DEM}}}{\ell_i^2},
\qquad
\Usrc_i=\bm0.
\]
Thus
\[
a_i\U_i=\bm0
\quad\Longrightarrow\quad
\U_i\approx\bm0.
\]
&
The solid-gas matrix coupling is removed, corresponding to
\[
\bm n\cdot D\nabla\U=0
\]
at the internal solid-gas interface. Gas cells are not hard-set to zero during the main solve, but they are zeroed in the final post-processing step when $\epsp_i\ge\epsp_c$.
\\

14. Boundary/interface condition
&
The zero face coefficient imposes an internal no-diffusion condition:
\[
\bm n\cdot\gamma\nabla\U=0
\]
between solid and non-solid cells.
&
The removed face coefficients impose
\[
\bm n\cdot D\nabla\U=0
\]
between the $\chi=1$ and $\chi=0$ regions. The code also removes the corresponding coupling on processor boundaries.
\\

15. Linear-solver repetitions
&
The equation is solved
\[
N_A=
\texttt{nUsInterpolationCorrectors}
\]
times. Before each solve, equation relaxation is applied.
&
The main hard-constrained Laplace system is solved up to five times. Let $\bm r_0^{(k)}$ be the initial residual vector at repetition $k$. The stopping test is
\[
\norm{\bm r_0^{(k)}}<10^{-3}
\]
and
\[
\left|
\frac{\norm{\bm r_0^{(k-1)}}}
{\max(\norm{\bm r_0^{(k)}},\VSMALL)}
-1
\right|<10^{-1}.
\]
\\

16. Optional corrector
&
No separate pseudo-transient equation is used. Repeated solves apply the same anchored steady equation (A).
&
When
\[
N_B=\texttt{nLaplaceSetValuesCorrectors}>0,
\]
the code solves
\[
c_t\frac{\partial\U}{\partial t}
-
\nabla\cdot(D\nabla\U)
=
\bm0,
\qquad
c_t=10^{-4}\,\mathrm{s},
\tag{B-corr}
\]
for $N_B$ repetitions. In a first-order finite-volume form:
\[
\frac{c_tV_i}{\Delta t}
\left(\U_i^{m+1}-\U_i^m\right)
+
\sum_j w_{ij}^{B}
\left(\U_i^{m+1}-\U_j^{m+1}\right)
=
\bm0.
\]
The \texttt{setValues} constraints are not re-applied inside this corrector system.
\\

17. Final overwrite
&
After the solve:
\[
\U_i\leftarrow
\begin{cases}
\Ud_i,&n_i>0.5,\\
\bm0,&n_i\le0.5\ \text{and}\ \epsp_i\ge\epsp_c,\\
\U_i,&\text{otherwise}.
\end{cases}
\tag{A-post}
\]
&
The same final overwrite is applied:
\[
\U_i\leftarrow
\begin{cases}
\Ud_i,&n_i>0.5,\\
\bm0,&n_i\le0.5\ \text{and}\ \epsp_i\ge\epsp_c,\\
\U_i,&\text{otherwise}.
\end{cases}
\tag{B-post}
\]
Thus occupied and gas-cell outputs are made exact after either method.
\\

18. Mathematical character
&
\textbf{Soft constrained screened Laplace:}
\[
\text{diffusion}+\text{reaction penalty}.
\]
The anchor strength is finite and mesh-scaled.
&
\textbf{Hard constrained Laplace:}
\[
\text{pure harmonic extension}
\]
with exact internal Dirichlet values imposed by replacing matrix rows.
\\

\bottomrule
\end{longtable}
\end{landscape}

\section{The essential mathematical difference}

The two equations can be written compactly as
\begin{align}
\texttt{laplaceAnchored:}\qquad
&
\mathcal L_A(\U)
=
-\nabla\cdot(\gamma\nabla\U)
+a(\U-\Usrc)
=
\bm0,
\label{eq:anchcompact}\\
\texttt{laplaceSetValues:}\qquad
&
\mathcal L_B(\U)
=
\nabla\cdot(D\nabla\U)
=
\bm0,
\qquad
\U_i=\Ud_i\quad(i\in\mathcal A).
\label{eq:setcompact}
\end{align}

Therefore:

\begin{itemize}[leftmargin=1.5em]
\item In \texttt{laplaceAnchored}, DEM information enters through a finite penalty
$a_i(\U_i-\Ud_i)$.
\item In \texttt{laplaceSetValues}, DEM information enters as an exact algebraic
constraint $\U_i=\Ud_i$.
\item In an empty solid region, \texttt{laplaceAnchored} contains a tiny
zero-directed reaction term, while \texttt{laplaceSetValues} is a pure Laplace equation.
\item For $c_{\mathrm{DEM}}\rightarrow\infty$ and
$c_{\mathrm{bg}}\rightarrow0$, the anchored model approaches the hard-constrained
Laplace model, provided that both methods use the same solid domain and the same
reference cells.
\end{itemize}

\section{One-dimensional interpretation}

Consider a one-dimensional empty solid interval $0<x<L$ between two known DEM
velocities $\U_L$ and $\U_R$.

\subsection*{\texttt{laplaceSetValues}}

With constant $D$,
\[
D\frac{\dd^2\U}{\dd x^2}=\bm0,
\qquad
\U(0)=\U_L,
\qquad
\U(L)=\U_R.
\]
Hence
\[
\boxed{
\U(x)
=
\U_L+\frac{x}{L}\left(\U_R-\U_L\right)
}
\]
which is exactly linear.

\subsection*{\texttt{laplaceAnchored}}

Inside an empty solid region with constant weak anchor $a_{\mathrm{bg}}$ and
$\Usrc=\bm0$,
\[
-\gamma\frac{\dd^2\U}{\dd x^2}
+a_{\mathrm{bg}}\U
=
\bm0.
\]
Define
\[
\lambda^2=\frac{a_{\mathrm{bg}}}{\gamma}.
\]
Then
\[
\U(x)
=
\bm C_1\cosh(\lambda x)
+
\bm C_2\sinh(\lambda x).
\]
When $a_{\mathrm{bg}}$ is extremely small,
$\lambda L\ll1$, so
\[
\cosh(\lambda x)\approx1,
\qquad
\sinh(\lambda x)\approx\lambda x,
\]
and the result becomes almost linear. It is not mathematically identical to
the pure Laplace result unless $a_{\mathrm{bg}}=0$ and the DEM anchors become exact.

\section{Conditions under which both methods become nearly equivalent}

The outputs approach each other when all of the following hold:
\begin{align}
c_{\mathrm{DEM}}&\rightarrow\infty,\\
c_{\mathrm{bg}}&\rightarrow0,\\
s_i&=\chi_i\quad\text{for every cell},\\
\{i:o_i=1\}&=\{i:\norm{\Ud_i}>0\},\\
N_B&=0,
\end{align}
and both methods use identical boundary conditions, mesh, numerical schemes,
linear solvers, tolerances, relaxation, and parallel decomposition.

The fourth condition fails for a stationary particle:
\[
o_i=1
\quad\text{but}\quad
\norm{\Ud_i}=0.
\]
In that case \texttt{laplaceAnchored} recognizes the cell as a DEM anchor, whereas
\texttt{laplaceSetValues} does not anchor it during the main solve.

\section{Compact comparison}

\begin{center}
\begin{tabularx}{0.96\textwidth}{L{4.5cm}YY}
\toprule
Property & \texttt{laplaceAnchored} & \texttt{laplaceSetValues}\\
\midrule
PDE type & Screened Laplace / reaction-diffusion & Pure Laplace\\
DEM constraint during main solve & Soft, finite penalty & Hard, exact row replacement\\
Empty-solid background bias & Tiny bias toward zero & No volumetric bias\\
Gas handling during solve & Strong zero anchor and disconnected diffusion & Interface disconnection; final gas reset\\
Stationary DEM particle & Recognized using $n_i$ & Missed by $\norm{\Ud_i}>0$ in main solve\\
Cell-size dependence & Anchor scales as $V_i^{-2/3}$ & Constant $D$; mesh enters through FV geometry\\
Optional pseudo-time step & No & Yes\\
Final occupied-cell value & Exactly $\Ud_i$ & Exactly $\Ud_i$\\
\bottomrule
\end{tabularx}
\end{center}

\section{Implementation note}

The model-selection line in the supplied file is written as
\[
\texttt{lookupOrDefault<word>("laplaceAnchored", "interpolationMode")}.
\]
For the usual OpenFOAM argument order \texttt{(keyword, defaultValue)}, the intended
form appears to be
\[
\boxed{
\texttt{lookupOrDefault<word>("interpolationMode", "laplaceAnchored")}
}
\]
Otherwise the code may look for a dictionary keyword named
\texttt{laplaceAnchored} and use \texttt{interpolationMode} as the default value.

\end{document}
